\section{Supplementary Problem 2.2}
Figure P2.2 in the reference document specifies
\begin{equation}
  x(n)=\begin{cases}
    1, & n\in\{-1,0,1,2\},\\
    \tfrac12, & n\in\{3,4\},\\
    0, & \text{otherwise}.
  \end{cases}
\end{equation}
The following results apply at integer time indices; all samples not
listed in a result are zero. The unit-step convention is $u(0)=1$.

\subsection{Part (a): Delay by two samples}
For $y_a(n)=x(n-2)$, every original time $m$ maps to $m+2$:
\begin{equation}
  y_a(n)=\begin{cases}
    1, & n\in\{1,2,3,4\},\\
    \tfrac12, & n\in\{5,6\},\\
    0, & \text{otherwise}.
  \end{cases}
\end{equation}

\subsection{Part (b): Folding followed by a four-sample delay}
For $y_b(n)=x(4-n)$, the sample mapping is $m\mapsto4-m$:
\begin{equation}
  y_b(n)=\begin{cases}
    \tfrac12, & n\in\{0,1\},\\
    1, & n\in\{2,3,4,5\},\\
    0, & \text{otherwise}.
  \end{cases}
\end{equation}

\subsection{Part (c): Advance by two samples}
For $y_c(n)=x(n+2)$, every original time $m$ maps to $m-2$:
\begin{equation}
  y_c(n)=\begin{cases}
    1, & n\in\{-3,-2,-1,0\},\\
    \tfrac12, & n\in\{1,2\},\\
    0, & \text{otherwise}.
  \end{cases}
\end{equation}

\subsection{Part (d): Multiplication by a reversed step}
Since $u(2-n)=1$ precisely when $n\leq2$, the product retains the
original samples up to and including $n=2$:
\begin{equation}
  y_d(n)=x(n)u(2-n)=\begin{cases}
    1, & n\in\{-1,0,1,2\},\\
    0, & \text{otherwise}.
  \end{cases}
\end{equation}
The samples of amplitude $1/2$ at $n=3$ and $n=4$ are removed.

\subsection{Part (e): Multiplication by a shifted impulse}
The impulse $\delta(n-3)$ selects only $n=3$. Evaluating the other
factor at this time gives $x(3-1)=x(2)=1$, so
\begin{equation}
  y_e(n)=x(n-1)\delta(n-3)=\delta(n-3).
\end{equation}
The output has a single sample of amplitude one at $n=3$.

\subsection{Part (f): A squared time argument}
The signal $y_f(n)=x(n^2)$ evaluates $x$ at the time index $n^2$;
it does not square the sample amplitudes. For integer $n$, the relevant
square indices are $0,1,4$:
\begin{equation}
  y_f(n)=\begin{cases}
    1, & n\in\{-1,0,1\},\\
    \tfrac12, & n\in\{-2,2\},\\
    0, & \text{otherwise}.
  \end{cases}
\end{equation}
In particular, $n=\pm2$ samples $x(4)=1/2$, and $|n|\geq3$ gives
an argument at least nine, where $x$ is zero.

\subsection{Parts (g) and (h): Even and odd components}
The required components are
\begin{align}
  x_e(n)&=\frac{x(n)+x(-n)}{2},\\
  x_o(n)&=\frac{x(n)-x(-n)}{2}.
\end{align}
They have the following sample values:
\begin{table}[H]
  \centering
  \caption{Even and odd components of the signal in Figure P2.2}
  \label{tab:p22-even-odd}
  \begin{tabular}{lrrrrrrrrr}
    \toprule
    $n$ & $-4$ & $-3$ & $-2$ & $-1$ & 0 & 1 & 2 & 3 & 4\\
    \midrule
    $x(n)$ & 0 & 0 & 0 & 1 & 1 & 1 & 1 & $1/2$ & $1/2$\\
    $x(-n)$ & $1/2$ & $1/2$ & 1 & 1 & 1 & 1 & 0 & 0 & 0\\
    $x_e(n)$ & $1/4$ & $1/4$ & $1/2$ & 1 & 1 & 1 & $1/2$ & $1/4$ & $1/4$\\
    $x_o(n)$ & $-1/4$ & $-1/4$ & $-1/2$ & 0 & 0 & 0 & $1/2$ & $1/4$ & $1/4$\\
    \bottomrule
  \end{tabular}
\end{table}
The even component satisfies $x_e(-n)=x_e(n)$, while the odd component
satisfies $x_o(-n)=-x_o(n)$ and $x_o(0)=0$. Adding the two components
recovers the original signal at every time index.

\subsection{Scilab implementation}
The script evaluates the given signal at the required time arguments,
uses elementwise products for parts (d) and (e), and uses
\texttt{n.\^{}2} to square each time index in part (f). Two figure
windows show parts (a)--(d) and (e)--(h), respectively. All plots use
the same time interval and amplitude limits. The script also prints
the values for parts (f)--(h) and the reconstruction error
$\max_n|x(n)-x_e(n)-x_o(n)|$, which is expected to be zero.
\lstinputlisting[caption={Scilab implementation of supplementary Problem 2.2},label={lst:p22}]{../2_2.sci}

\subsection{Execution evidence}
\resultfigure{images/problem22-ad.png}
  {Time shifts, folding, and step multiplication in Problem 2.2(a)--(d)}
  {fig:p22-ad}
\resultfigure{images/problem22-eh.png}
  {Impulse multiplication, squared time argument, and even and odd components in Problem 2.2(e)--(h)}
  {fig:p22-eh}
The captured plots agree with the analytical results for all eight parts.
In particular, part (e) contains one impulse at $n=3$, part (f) has
amplitudes $1/2,1,1,1,1/2$ at $n=-2,-1,0,1,2$, and the even and odd
components show the required symmetry. The numerical reconstruction
error is a separate console check and is not displayed in these figures.
